%% kapa: exam=1 round=31 year=2020 period=2 form=A subject=appraisal_law subject_ko=감정평가관계법규 seq=14 number=14
%% source: Q-Net 감정평가사 기출문제 (cst00309) artlSeq=5207219 file=extracted/1008/5207219/A형/A형/2020년도 제31회 감정평가사 1차 시험_2교시_A형.pdf
%% conversion: kordoc 4.24.51 -> tools/gen1.py (auto); review=proofread-sonnet; flags=none
\begin{question}{14}
감정평가 및 감정평가사에 관한 법령상 감정평가에 관한 설명으로 옳지 않은 것은?
\begin{choices}
  \item 감정평가업자가 토지를 감정평가하는 경우 적정한 실거래가가 있는 경우에는 이를 기준으로 할 수 있다.
  \item 감정평가업자는 해산 또는 폐업하는 경우에도 감정평가서 관련 서류를 발급일부터 5년 이상 보존하여야 한다.
  \item 국토교통부장관은 감정평가 타당성조사를 할 경우 해당 감정평가를 의뢰한 자에게 의견진술기회를 주어야 한다.
  \item 감정평가법인은 감정평가서를 의뢰인에게 발급하기 전에 같은 법인 소속의 다른 감정평가사에게 감정평가서의 적정성을 심사하게 하여야 한다.
  \item 토지 및 건물의 가격에 관한 정보 및 자료는 감정평가 정보체계의 관리대상에 해당한다.
\end{choices}
\end{question}
\answer{2}
